\subsection{Cuadros incluidos, contextuales y descartados}

El diccionario actualizado codifica el papel analítico de los cuadros. 1HM, 2HM, 5PM, 6V\_MPIO, 8V\_MPIO y 10V\_MPIO generan candidatas; 17PM queda como medida educativa descriptiva y candidata reservada para sensibilidad. 1PM se usa como contexto y fuente de denominadores, y 1V\_MPIO únicamente como contexto. 6PM se excluye: sus variables son equivalentes a 5PM pero organizadas por grupos de edad, por lo que incluirlo como otra medida duplicaría la dimensión observada. El IPM se guarda para validación externa; sus 15 privaciones no se suman al PCA.

En 6V\_MPIO se mantiene el conjunto de seis categorías acordado por el equipo: guadua; tapia pisada/bahareque/adobe; madera burda/tabla/tablón; caña, esterilla y otros vegetales; materiales de desecho; y sin paredes. La elección agrupa materiales que requieren interpretación conjunta y no afirma que tengan idéntica severidad. El solapamiento directo con la privación IPM de paredes se evaluará externamente.

En 8V\_MPIO se analizan por separado energía, gas e internet. Acueducto y alcantarillado quedan fuera de la matriz principal por cercanía conceptual con privaciones IPM, pero se conservan como variables descriptivas. Las categorías \enquote{Sin información} permanecen registradas; no se confunden con viviendas que carecen del servicio.

La decisión confirmada el 18 de octubre asigna recolección de basuras únicamente a sensibilidad exploratoria. La matriz histórica de siete variables la incluía; las ocho actuales la retiran e incorporan 2HM y 6V, además de revisar fórmulas de actividad y estrato. Por ello, los diagnósticos históricos no son aplicables automáticamente a la especificación actual. Este acuerdo posterior se registra en la actualización del informe y en su nota de decisiones; el CSV histórico se conserva para trazabilidad.

\subsection{Dirección conceptual}

Las candidatas se orientan para que valores altos indiquen mayor condición potencial de vulnerabilidad según la operacionalización acordada. Esta orientación es una convención explícita del estudio, no una constatación causal. En particular, tener menores o personas mayores en el hogar representa una posible carga de cuidado, pero no demuestra dependencia económica ni funcional. El estrato es una clasificación administrativa del inmueble residencial asociada con servicios públicos; la inclusión de \enquote{no tiene estrato} como vulnerabilidad es una decisión del proyecto, no una consecuencia automática de la Ley 142 ni una medición directa del ingreso \citep{dane_estratificacion}.

En 5PM, la variable censal clasifica la actividad realizada la semana anterior por población de 10 años y más. El agregado propuesto mezcla trabajo no remunerado, búsqueda de empleo e incapacidad permanente con mayor ponderación de esta última. Por tanto, es una puntuación experimental municipal y no debe nombrarse tasa de desempleo, dependencia económica del hogar ni \enquote{última actividad laboral del hogar}. El cuadro no contiene el vínculo persona-hogar necesario para esa última interpretación \citep{dane_5pm}.

17PM mide nivel educativo alcanzado en población de cinco años y más. Sus categorías por municipio, área y sexo no están desagregadas por edad en el cuadro utilizado. La proporción agregada puede variar por la composición etaria además de las oportunidades educativas; por ello no se interpreta como educación ajustada por edad ni como asistencia escolar. La candidata agregada \texttt{sin\_media\_completa\_17PM} queda fuera del primer PCA, pendiente de análisis exploratorio de sensibilidad.
